\begin{table}[htbp]
\centering
\small
\caption{Двадцать шесть пунктов ограничений, свёрнутых в семь групп основного текста}
\label{tab:s16}
\begin{tabular}{>{\RaggedRight\hspace{0pt}\arraybackslash}p{0.302\linewidth}>{\RaggedRight\hspace{0pt}\arraybackslash}p{0.302\linewidth}>{\RaggedRight\hspace{0pt}\arraybackslash}p{0.302\linewidth}}
\toprule
Группа & Какие пункты сворачивает & Что она говорит рецензенту \\
\midrule
1. Идентифицируемость дизайна & 2, 12, 13 & два первичных исхода из трёх неидентифицируемы, негативный контроль невыполним, пять holdout одноклассовые и три из них совпадают \\
2. Содержательная валидность признаков & 1, 17 & индекс покрывает две категории из тринадцати полностью, четырнадцать признаков не посчитаны \\
3. Зависимость от процедуры & 3, 4, 6, 7, 9, 10 & процедуры 2–4 exploratory, 1 и 2 — одно семейство, net есть не у всех, эффекты малы, поправка Бонферрони, разные популяции контрастов \\
4. Источник и жанр & 11, 14, 15, 19, 21 & дефект извлечения односторонний, пропуски несут метку источника, четыре машинных кластера, разметка уровня по источнику одним человеком, человеческая часть до 2022 года \\
5. Ограниченная воспроизводимость судьи & 8 & ICC(2,1) 0.477, пятнадцать значений медианы шкалы, 85.9\% в двух точках \\
6. Отсутствие проспективного bridge set & 5 & все процедуры на одном корпусе: альтернативные измерения, не независимые выборки \\
7. Границы стресс-теста & 22, 23, 24, 25, 26 & дискретизация терцилей, нормирование форматных счётчиков на объём, знаменатель 10 из 16, составные преобразования объёма, отклонённая на preflight sensitivity-проверка длины \\
\bottomrule
\end{tabular}
\begin{minipage}{\linewidth}\footnotesize\vspace{2pt}
\textit{Примечание.} Единица анализа — группа ограничений. Знаменатель: 26 пунктов полного списка, семь групп в основном тексте. Данные: список ограничений от 2026-07-29 с дополнениями 31 июля и 1 августа. Таблица описывает состав; чисел прогона в ней нет. Шкала: шкалы нет. Сокращения: ICC(2,1) — коэффициент внутриклассовой корреляции для случайных оценщиков; bridge set — независимая выборка для переноса выводов. Пропуски: пункты 16, 18 и 20 в основной текст не вошли и остаются только в полном списке. Поправка на множественные сравнения не применялась.
\end{minipage}
\end{table}
